
Execution-feedback reinforcement learning and direct preference optimization are both used to post-train code generation models, yet existing evaluation practice relies on pass@1 benchmark outcomes, which cannot distinguish whether a given alignment method induces structurally different policy change or merely increases confidence about existing surface patterns. This paper introduces \textit{structural efficiency of policy movement} — defined as semantic AST edit distance per unit KL divergence — as a diagnostic metric for code generation alignment, and describes an end-to-end measurement framework combining FA-AST node classification, ZSS tree edit distance, KL-matched checkpoint comparison, and bootstrap statistical testing. The framework is implemented and validated on DeepSeek-Coder-7B-instruct-v1.5 using TRL GRPOTrainer and DPOTrainer. Two sub-experiments were executed. The first (h-e1) demonstrated that the measurement pipeline executes correctly end-to-end on real checkpoints, producing a 95\% bootstrap CI of [4.65, 8.73] for the mean edit-per-KL differential on proof-of-concept data; however, this experiment used hand-crafted code completions with engineered structural differences and does not constitute empirical evidence about real GRPO training behavior. The second (h-m1) measured Semantic Edit Proportion (SEP) — the fraction of edits targeting control-flow and data-flow AST nodes — on real checkpoint outputs, yielding SEP ≈ 0.237 for both GRPO and DPO (Mann-Whitney U = 18,346.5, p = 0.4248); this result is underpowered due to checkpoint aliasing (effective n ≈ 2 from 27 nominal pairs). The paper reports checkpoint aliasing as a previously undocumented confound in checkpoint-comparison studies of RL fine-tuning and provides a corrected experimental protocol. The core hypothesis that execution reward selectively concentrates policy movement on semantically relevant AST nodes remains empirically untested; the measurement infrastructure is validated and ready for application to a corrected experimental run.
